%% ECON 282E -- Homework 1 mini-paper template. At most six pages, references included.
%% Every number must be \input from output/, written by run_all.py -- never typed by hand.
\documentclass[11pt]{article}
\usepackage[margin=1in]{geometry}
\usepackage{amsmath,amssymb,booktabs,graphicx}
\usepackage[T1]{fontenc}
\usepackage{microtype}
\usepackage[colorlinks=true,linkcolor=blue!50!black,urlcolor=blue!50!black,citecolor=blue!50!black]{hyperref}

\title{Title: A Question, Stated as a Claim}
\author{Your Name\thanks{ECON 282E, UC Riverside, Fall 2026. Code and data: \texttt{run\_all.py} reproduces every number.}}
\date{\today}

\begin{document}
\maketitle

\begin{abstract}
Three sentences: the question, what you did, and the answer --- with the number.
\end{abstract}

\section{Question}
Why the question matters, whether it is positive or normative, and the one number that answers it.

\section{Model and Equilibrium}
The environment, the equilibrium definition, and what the model assumes away.

\section{Calibration}
% \input{output/table_calibration.tex}
Targets, sources (series ID, sample, download date), and the choices that move the results.

\section{Solution Method and Its Accuracy}
% \input{output/table_methods.tex}
Which method, why it is fit for each question, and the evidence that the solution is right.

\section{Results}
% \input{output/table_answers.tex}
The positive answer, the normative answer, and how sensitive each is to the choices above.

\section{What the Model Cannot Say}
The limits of the representative-agent answer, and what model you would need next.

\bibliographystyle{apalike}
% \bibliography{references}
\end{document}
